\subsection{eBPF as the Low-Level Observability Substrate}
\label{subsec:ebpf-foundation}

The fundamental limitation of application-level observability (OSI Layer 7) resides in its inherent detachment from the underlying operating system infrastructure. To definitively close the ``Layer Visibility Gap'' identified in the PANOPTES architecture, the framework incorporates the Extended Berkeley Packet Filter (eBPF) as its core low-level instrumentation substrate. eBPF is a revolutionary kernel technology that permits the execution of custom, sandboxed bytecode directly within the Linux kernel, dynamically altering the operating system's behavior without requiring source code modifications or risky kernel modules \cite{rice2023learning}.

In cloud-native ecosystems, containers are essentially isolated Linux processes sharing a single host kernel \cite{salazar2022security}. Consequently, an eBPF program attached to kernel events inherently possesses complete, unfiltered visibility over all containerized workloads operating on that node \cite{rice2023learning, salazar2022security}. To guarantee system stability and security, eBPF relies on an in-kernel verifier that performs rigorous static analysis on the bytecode, ensuring memory safety, preventing infinite loops, and strictly validating pointer dereferences before execution \cite{calavera2019linux}. Once verified, the bytecode is translated by a Just-In-Time (JIT) compiler into native machine instructions, enabling extremely high-performance telemetry collection without the prohibitive overhead of context-switching between user and kernel space \cite{rice2023learning, calavera2019linux}.

The strategic advantage of eBPF in the PANOPTES architecture lies in its ability to bridge the gap between network mechanics and application semantics. Traditional security and performance monitoring tools struggle with the ephemeral nature of Kubernetes pods and NAT (Network Address Translation) obfuscation \cite{salazar2022security}. eBPF overcomes this by tracking network connections directly at the socket layer (Layer 4), extracting high-fidelity metadata (such as TCP retransmissions, connection state, and RTT) pre-NAT, and seamlessly mapping these low-level network events back to Kubernetes identities, including pod names and namespaces \cite{salazar2022security, shekhar2023ultimate}.

By hooking into specific kernel functions via \textit{kprobes}, \textit{tracepoints}, and socket-level events, eBPF acts as an unbroken chain of evidence from the network layer up to the application layer \cite{calavera2019linux}. In the context of PANOPTES, integrating eBPF network telemetry guarantees that structural Layer 4 network degradations---such as the $2000\,\text{ms}$ socket latency injected via Chaos Mesh---are deterministically captured and correlated with Layer 7 distributed traces \cite{shekhar2023ultimate}. This unified, cross-layer visibility transforms the network from an opaque component into a mathematically observable medium, effectively resolving the architectural blindness observed in pure L7 instrumentation.